% !TEX root = main.tex

\begin{abstract}
Railway infrastructure faces risks from rail defects, potentially leading to rail breaks and reduced efficiency. This paper presents a probabilistic forecasting model using a covariate-dependent Markov jump process to predict defect severity propagation. We explore different generator parameterizations, each influencing model flexibility and prediction accuracy. The models are evaluated using probabilistic scoring rules (log score, Brier score, and ranked probability score) and benchmarked against several reference models, including cumulative link models. The Markov jump process achieves the highest prediction accuracy, with further improvements through ensemble forecasting. We demonstrate that incorporating spatial heterogeneity, such as local rail conditions, consistently improves forecasting accuracy across all model settings. Additionally, we conduct a thorough diagnostic analysis of forecast quality, focusing on calibration and sharpness. This study offers a robust framework for predictive maintenance with the potential to improve both track safety and operational efficiency.
\end{abstract}